\section{Introduction}

Quantum chromodynamics (QCD) is the theory of the
strong nuclear force that connects hadronic bound states to their partonic 
constituents, quarks and gluons. Although quarks and gluons cannot be 
directly accessed in experiments, the connection between hadrons 
and partons can be characterized through parton 
distribution functions (PDFs).

PDFs capture aspects of hadron structure associated with the momentum, 
angular momentum and spin of the constituent quarks and gluons, and play a 
central role in our understanding of high energy hadronic scattering processes  
(see, for example, 
\cite{Brock:1993sz,Sterman:1994ce,Collins:2011zzd}). Through factorization, the 
scattering 
amplitudes of simple scattering processes, such as deep inelastic scattering 
and Drell-Yan production, can be expressed as the convolution of perturbative 
coefficients and PDFs, which encapsulate the nonperturbative dynamics of QCD at 
hadronic scales.

In principle, the direct calculation of PDFs from QCD will provide new insight 
into hadronic structure, more stringent tests of QCD, and reduced systematic 
uncertainties in high energy scattering experiments. At present, however, the 
only systematic method for ab initio, nonperturbative QCD calculations is 
lattice QCD, in which QCD is formulated on a discrete Euclidean hypercube. PDFs 
are defined in terms of matrix elements of light-front wave functions and cannot 
be directly determined from Euclidean lattice QCD. PDFs are currently
determined from global analyses of a wide range of scattering data (see, for 
example, 
\cite{Arbabifar:2013tma,Ball:2013lla,Jimenez-Delgado:2013boa,
Jimenez-Delgado:2014xza,Dulat:2015mca, 
Hou:2016sho,Shahri:2016uzl} for a selection of recent analyses).

Lattice QCD calculations have instead focused on the first 
few Mellin moments of PDFs, which can be related to matrix elements of local 
twist-two operators, where twist is the dimension minus the spin of the 
operator. The lattice 
regulator breaks rotational symmetry, which induces mixing between lattice 
operators that would not mix in the continuum. The mixing between twist-two 
operators of different 
mass dimension introduces power-divergent mixing on the lattice, preventing the 
extraction of more than three moments of PDFs 
\cite{Detmold:2001dv,Detmold:2003rq}.

Recently a new approach to determining PDFs on the lattice was proposed, via 
Euclidean counterparts of PDFs generally referred to as quasi PDFs  
\cite{Ji:2013dva,Ji:2014gla,Lin:2014zya,Gamberg:2014zwa,Ji:2015jwa,
Alexandrou:2015rja,Chen:2016utp}. A similar framework was proposed in 
\cite{Ma:2014jla}. The 
quasi PDFs are Euclidean matrix elements determined at finite nucleon 
momentum. At large Euclidean momentum, the quasi PDFs can be related to the 
true PDFs through an effective theory expansion, the Large Momentum Effective 
Theory (LaMET).

Preliminary lattice calculations have been encouraging 
\cite{Lin:2014zya,Alexandrou:2015rja}, although both 
calculations have incorporated only a single lattice spacing and a full 
understanding of systematic uncertainties is far from complete. In particular, 
there are three challenges for the approach as it stands: the 
restriction to low nucleon momentum with the computational resources currently 
available; a full understanding of the renormalisation of extended Euclidean 
operators; and the precise relation between light-front PDFs and Euclidean quasi 
PDFs. 

These difficulties can be broadly classified as either chiefly practical, or 
chiefly theoretical. The first challenge, that associated with the systematic 
uncertainties corresponding to low values of nucleon momentum on current 
lattices, is largely a practical issue. Studies in the spectator 
di-quark model \cite{Gamberg:2014zwa} suggest that moderate improvements in 
computational resources, and new algorithms tailored to nucleons with large 
momentum \cite{Bali:2016lva}, 
will likely solve this difficulty, at least to a precision that can contribute 
to global analyses of the PDFs in regions of parameter space that are 
experimentally inaccessible. We will not consider these practical 
difficulties any further and focus instead on the theoretical 
aspects of quasi PDFs.

We address one of the theoretical challenges by proposing a new approach to 
calculating quasi PDFs on the lattice, in which the lattice degrees of freedom 
are smeared via the gradient flow 
\cite{Narayanan:2006rf,Luscher:2011bx,Luscher:2013cpa}. Using ringed 
fermions, which do not require any multiplicative wavefunction renormalization
\cite{Makino:2014taa,Hieda:2016lly}, the corresponding lattice matrix elements 
remain finite in the continuum limit. This approach evades the problem of the 
power-divergence associated with the Wilson line operator that defines the 
quasi PDF. The renormalization of quasi PDFs has been viewed through the lens 
of heavy quark effective theory \cite{Ji:2015jwa} and, more recently, a 
counterterm procedure has been proposed to remove this power-divergence 
\cite{Chen:2016fxx,Ishikawa:2016znu}, but neither approach has been established 
beyond two loops in perturbation theory.


Here we examine the relation between the smeared quasi 
PDF and the light-front PDF, and focus on the limit in which the flow 
time is small compared to the 
length scale set by the nucleon momentum and the nucleon momentum is 
sufficiently large that higher twist effects can be neglected. In this limit, 
we express the moments of the smeared quasi PDF in terms of moments of the 
light-front PDF via a small flow-time expansion.
The primary advantage of 
our approach is the finite continuum limit for nonperturbative matrix elements 
determined using lattice QCD. The matching between the smeared quasi PDF, 
regulated by the flow time, and the light-front PDF in, say, the $\ms$ 
scheme can be carried out in the continuum and is independent of the details of 
the nonperturbative lattice calculation.

We start by revisiting the definitions of the light-front and quasi PDFs in 
Section \ref{sec:defs}. We then analyze the relation between the 
light-front and quasi PDFs in Section \ref{sec:spdf} and derive an evolution 
equation for the matching kernel in Section \ref{sec:match}. We 
present our summary and conclusions in Section \ref{sec:concs}.
